\documentclass[10.5pt,a4paper]{article}
\usepackage{fontspec}
\setmainfont{Arial}
\usepackage[margin=1.2cm,top=1.2cm,bottom=1.2cm]{geometry}
\usepackage{listings}

\setlength{\parindent}{0pt}
\setlength{\parskip}{3pt}

\lstset{
    basicstyle=\ttfamily\small,
    frame=single,
    breaklines=true,
    showstringspaces=false,
    tabsize=4
}

\begin{document}

% ==================== SƏHİFƏ 1 ====================
\begin{center}
    {\large \textbf{BAKI ALİ NEFT MƏKTƏBİNİN NƏZDİNDƏ}}\\[0.1cm]
    {\LARGE \textbf{RƏQƏMSAL BİLİKLƏR LİSEYİ}}\\[0.2cm]
    {\textbf{Fənn:} İnformatika (C++ Proqramlaşdırma) \quad | \quad \textbf{Sinif:} 8-ci sinif}\\[0.05cm]
    {\textbf{Mövzu:} Dərs 01 -- C++ Sintaksisinin Əsasları və Xətti Alqoritmlər}\\[0.15cm]
    {\Large \textbf{EV TAPŞIRIĞI İŞ VƏRƏQİ -- № 01}}
\end{center}

\vspace{0.05cm}
\noindent\fbox{%
\parbox{0.98\textwidth}{%
\vspace{0.05cm}
\textbf{Şagirdin Adı, Soyadı:} \hrulefill \quad 
\textbf{Sinif:} \underline{\hspace{1.5cm}} \quad 
\textbf{Tarix:} \underline{\hspace{2cm}} \quad
\textbf{Qiymət:} \underline{\hspace{1.2cm}} / 100
\vspace{0.05cm}
}}

\vspace{0.2cm}
{\large \textbf{HİSSƏ A: Nəzəriyyə və Kod Analizi (30 Bal)}}
\vspace{0.05cm}
\hrule
\vspace{0.15cm}

\textbf{1. C++ Proqramının Əsas Quruluşu (5 bal)}\\
C++ proqramının başında yazılan \texttt{\#include <iostream>} kitabxanası və \texttt{using namespace std;} təyinatı nə üçündür? Qısa izah edin.\\
\noindent\rule{\textwidth}{0.4pt}\\[0.15cm]
\noindent\rule{\textwidth}{0.4pt}

\vspace{0.2cm}
\textbf{2. Xüsusi Simvollar (Escape Sequences) (5 bal)}\\
Aşağıdakı simvolların proqram çıxışında nə rol oynadığını qarşısına qeyd edin:
\begin{itemize}
    \item \texttt{\textbackslash n} simvolu: \hrulefill
    \item \texttt{\textbackslash t} simvolu: \hrulefill
    \item \texttt{\textbackslash "} simvolu: \hrulefill
\end{itemize}

\vspace{0.15cm}
\textbf{3. Konsol Çıxışını Əvvəlcədən Təyin Edin (Predict Output) (10 bal)}\\
Aşağıdakı kod parçası icra olunduqda konsolda hansı nəticələr görünəcək? Dəqiq qeyd edin:

\begin{minipage}[t]{0.52\textwidth}
\begin{lstlisting}[language=C++]
int a = 23;
int b = 4;
cout << "Netice 1: " << a / b << endl;
cout << "Netice 2: " << a % b << endl;
cout << "Netice 3: " << a + b * 2 << endl;
\end{lstlisting}
\end{minipage}
\hfill
\begin{minipage}[t]{0.44\textwidth}
\noindent\fbox{%
\parbox{0.95\linewidth}{%
\textbf{Konsol Çıxışı:}\\[0.1cm]
1. \hrulefill\\[0.3cm]
2. \hrulefill\\[0.3cm]
3. \hrulefill\\[0.3cm]
}}
\end{minipage}

\vspace{0.2cm}
\textbf{4. Sintaksis Xətalarını Aşkarlayın və Düzəldin (10 bal)}\\
Aşağıdakı kod parçasında 4 kobud xəta var. Hər xətanı göstərin və qarşısına düzgün formasını yazın:

\begin{minipage}[t]{0.52\textwidth}
\begin{lstlisting}[language=C++]
#include <iostream>
using namespace std;

int main() {
    int qiymet = 12.5;
    cin << say;
    cout >> "Cemi: " >> qiymet * say
    return 0
}
\end{lstlisting}
\end{minipage}
\hfill
\begin{minipage}[t]{0.44\textwidth}
\textbf{Düzəlişləriniz:}\\[0.1cm]
1. \hrulefill\\[0.3cm]
2. \hrulefill\\[0.3cm]
3. \hrulefill\\[0.3cm]
4. \hrulefill\\[0.3cm]
\end{minipage}

\vfill
\begin{center}
\footnotesize \textit{Müəllim: Avaz Əsgərov | BANM Nəzdində Rəqəmsal Biliklər Liseyi -- Səhifə 1 / 2}
\end{center}

\newpage

% ==================== SƏHİFƏ 2 ====================
{\large \textbf{HİSSƏ B: Praktik Proqramlaşdırma Məsələləri (70 Bal)}}
\vspace{0.05cm}
\hrule
\vspace{0.15cm}

\textbf{Məsələ 1: Valyuta Dəyişmə Məntəqəsi Kalkulyatoru (20 bal)}\\
İstifadəçidən manat məbləğini (\texttt{double aznAmount}) və dollar məzənnəsini (\texttt{double exchangeRate}) qəbul edən, neçə ABŞ dolları alındığını hesablayıb (\texttt{usd = azn / exchangeRate}) ekrana çıxaran C++ kodunu yazın:

\noindent\fbox{%
\parbox{0.98\textwidth}{%
\vspace{0.05cm}
\textbf{Məsələ 1 üçün C++ Kodu:}\\[4.6cm]
}}

\vspace{0.25cm}
\textbf{Məsələ 2: Avtomobil Səyahəti və Yanacaq Xərci (25 bal)}\\
Məsafəni (\texttt{km}), 100 km-ə sərfiyyatı (\texttt{litr}) və 1 litr benzinin qiymətini (\texttt{AZN}) daxil edərək, cəmi tələb olunan benzini (\texttt{(km / 100.0) * litr}) və ümumi xərci hesablayıb ekrana çıxaran proqram kodunu yazın:

\noindent\fbox{%
\parbox{0.98\textwidth}{%
\vspace{0.05cm}
\textbf{Məsələ 2 üçün C++ Kodu:}\\[4.6cm]
}}

\vspace{0.25cm}
\textbf{Məsələ 3: Saniyə Ölçən və Vaxt Bölgüsü (Böyük Qalıq Məsələsi) (25 bal)}\\
İstifadəçidən saniyə sayını (\texttt{int totalSeconds}) daxil etməsini istəyin. Tam bölmə (\texttt{/}) və qalıq (\texttt{\%}) ilə bu vaxtı \textbf{saat, dəqiqə və saniyə} formatında hesablayıb ekrana çıxaran C++ kodunu yazın:

\noindent\fbox{%
\parbox{0.98\textwidth}{%
\vspace{0.05cm}
\textbf{Məsələ 3 üçün C++ Kodu:}\\[4.6cm]
}}

\vfill
\begin{center}
\footnotesize \textit{Müəllim: Avaz Əsgərov | BANM Nəzdində Rəqəmsal Biliklər Liseyi -- Səhifə 2 / 2}
\end{center}

\end{document}
